\documentclass[11pt,a4paper]{article}
\usepackage[utf8]{inputenc}
\usepackage[T1]{fontenc}
\usepackage[spanish,es-minimal,es-noshorthands]{babel}
\usepackage{lmodern}
\usepackage[sfdefault,lining]{FiraSans}
\usepackage[a4paper,top=2.3cm,bottom=2.2cm,left=1.8cm,right=1.8cm,headheight=1.3cm,headsep=0.4cm,footskip=1cm]{geometry}
\usepackage{amsmath,amssymb}
\usepackage{icomma}
\usepackage{graphicx}
\usepackage[table]{xcolor}
\usepackage{tikz}
\usepackage{fontawesome5}
\usepackage{ccicons}
\usepackage{fancyhdr}
\usepackage{lastpage}
\usepackage{tasks}
\usepackage{needspace}
\usepackage[most]{tcolorbox}
\usepackage{array}
\usepackage[hidelinks]{hyperref}
\AtBeginDocument{\hypersetup{pdftitle={Ficha de ejercicios de la unidad 9: Perímetros y áreas. Teorema de Pitágoras. 3º ESO},pdfsubject={Matemáticas 3º ESO},pdfkeywords={teorema de Pitágoras, hipotenusa, catetos, perímetro, área, polígonos regulares, apotema, circunferencia, círculo, sector circular, corona circular}}}

\definecolor{navy}{HTML}{1B2748}
\definecolor{accent}{HTML}{E8603A}
\definecolor{teal}{HTML}{1F9E90}
\definecolor{soft}{HTML}{EEF1F8}
\definecolor{muted}{HTML}{5D6685}

\setlength{\parindent}{0pt}
\setlength{\parskip}{3pt}
\settasks{label=\alph*), label-format=\bfseries\color{navy}, label-width=1.3em, item-indent=1.8em, column-sep=0.6em, after-item-skip=5pt, before-skip=3pt, after-skip=4pt}

% ---------- capçalera i peu ----------
\newcommand{\llicencia}{\href{https://creativecommons.org/licenses/by-nc-sa/4.0/deed.es}{\ccbyncsa}\ \ Licencia Creative Commons Atribución-NoComercial-CompartirIgual 4.0 Internacional (CC BY-NC-SA 4.0)}
\pagestyle{fancy}
\fancyhf{}
\fancyhead[L]{\raisebox{-0.35\height}{\includegraphics[height=0.95cm]{logo-gva.png}}}
\fancyhead[C]{\small\color{muted}Ficha de ejercicios · Unidad 9: Perímetros y áreas. Teorema de Pitágoras · 3\textsuperscript{o} ESO}
\fancyhead[R]{\raisebox{-0.35\height}{\includegraphics[height=0.95cm]{logo-vermellar.png}}}
\fancyfoot[L]{\scriptsize\color{muted}\llicencia}
\fancyfoot[R]{\small\color{muted}\thepage}
\renewcommand{\headrulewidth}{0.4pt}
\renewcommand{\headrule}{\color{muted!40}\hrule height \headrulewidth}
\fancypagestyle{primera}{\fancyhf{}\renewcommand{\headrulewidth}{0pt}
  \fancyfoot[L]{\scriptsize\color{muted}\llicencia}
  \fancyfoot[R]{\small\color{muted}\thepage}}

% ---------- nivells i exercicis ----------
\newcommand{\estrelles}[1]{{\color{accent}\ifcase#1\or\faStar\faStar[regular]\faStar[regular]\or\faStar\faStar\faStar[regular]\or\faStar\faStar\faStar\fi}}
\newcounter{exer}
\newcommand{\ex}[1]{\par\needspace{4\baselineskip}\medskip\stepcounter{exer}\noindent{\color{navy}\bfseries\theexer.}\ \ {\small\estrelles{#1}}\ \ }

\newcommand{\seccio}[2]{\par\needspace{10\baselineskip}\bigskip
  \begin{tcolorbox}[colback=navy,colframe=navy,coltext=white,arc=2mm,boxrule=0pt,left=3mm,right=3mm,top=1.2mm,bottom=1.2mm]
  {\large\bfseries #1\quad #2}
  \end{tcolorbox}\vspace{-2pt}}

\newcolumntype{C}[1]{>{\centering\arraybackslash}p{#1}}
\usepackage{pifont}

% quadre de recordatori (només a la fitxa d'exercicis)
\newcommand{\recorda}[1]{\par\smallskip\begin{tcolorbox}[colback=accent!7,colframe=accent!55,arc=2mm,boxrule=0.5pt,left=2.5mm,right=2.5mm,top=1mm,bottom=1mm]\small{\color{accent}\faLightbulb[regular]}\ \textbf{Recuerda:}\ #1\end{tcolorbox}}
\newcommand{\res}[1]{#1}
\newcommand{\rest}[1]{#1}

\def\fmneed{14}
% =====================================================================
% Capçalera FentMates de les fitxes: un color per a cada curs.
% Els colors són els de la pàgina inicial (LEVELS a assets/js/materials.js);
% el logo ve de _eines/logo/pdf/ (python3 _eines/logo/genera.py).
%
% Ús, al final del preàmbul de la unitat:
%   % =====================================================================
% Capçalera FentMates de les fitxes: un color per a cada curs.
% Els colors són els de la pàgina inicial (LEVELS a assets/js/materials.js);
% el logo ve de _eines/logo/pdf/ (python3 _eines/logo/genera.py).
%
% Ús, al final del preàmbul de la unitat:
%   % =====================================================================
% Capçalera FentMates de les fitxes: un color per a cada curs.
% Els colors són els de la pàgina inicial (LEVELS a assets/js/materials.js);
% el logo ve de _eines/logo/pdf/ (python3 _eines/logo/genera.py).
%
% Ús, al final del preàmbul de la unitat:
%   \input{../comu/fentmates.tex}
%   \fmcurs{1eso}
% i, al començament del document:
%   \capcalera{Matemàtiques}{1r ESO}{Fitxa d'exercicis · Unitat 2}{Divisibilitat}
% Opcional abans de \input: \def\fmneed{14} (espai mínim abans de cada secció).
% =====================================================================
\makeatletter
\def\fm@logodir{../../logo/pdf}
\providecommand{\fmneed}{10}
\def\fm@cursnom{}

% clau del curs, color i color del text damunt del color
\newcommand{\fm@curs}[3]{\@namedef{fm@color@#1}{#2}\@namedef{fm@text@#1}{#3}}
\fm@curs{1eso}{E6B84A}{navy}
\fm@curs{2eso}{3B82F6}{white}
\fm@curs{3eso}{FF7A50}{white}
\fm@curs{4eso-a}{A78BFA}{white}
\fm@curs{4eso-b}{6D28D9}{white}
\fm@curs{1batx-mi}{9B1C1C}{white}
\fm@curs{1batx-ccss}{EF6B6B}{white}
\fm@curs{2batx-mii}{1E6B47}{white}
\fm@curs{2batx-ccss}{2BB3A3}{white}

\newcommand{\fmcurs}[1]{%
  \@ifundefined{fm@color@#1}{\PackageError{fentmates}{Curs desconegut: #1}{}}{}%
  \edef\fm@tmp{\noexpand\definecolor{curs}{HTML}{\@nameuse{fm@color@#1}}%
               \noexpand\colorlet{curstext}{\@nameuse{fm@text@#1}}}\fm@tmp
  \colorlet{cursfosc}{curs!72!black}%
  \def\fm@logo{\fm@logodir/vertical-#1.pdf}%
  \def\fm@marca{\fm@logodir/marca-#1.pdf}}

% ---------- capçalera de la primera pàgina ----------
% #1 matèria (petit), #2 curs, #3 tipus de document i unitat, #4 títol
\newcommand{\capcalera}[4]{%
  \gdef\fm@cursnom{#2}%
  \par\noindent
  \begin{tcolorbox}[colback=curs,colframe=curs,boxrule=0pt,arc=2.2mm,width=3.5cm,height=2.15cm,
      valign=center,left=3mm,right=2mm,top=1mm,bottom=1mm,nobeforeafter,box align=center]
    \color{curstext}{\footnotesize #1}\par\vspace{1pt}{\fontsize{19}{21}\selectfont\bfseries #2\par}%
  \end{tcolorbox}\hspace{5mm}%
  \begin{minipage}[c]{\dimexpr\linewidth-3.5cm-5mm-2.9cm\relax}
    {\color{cursfosc}\small\bfseries\MakeUppercase{#3}}\par\vspace{2pt}
    {\color{navy}\fontsize{20}{23}\selectfont\bfseries #4\par}
  \end{minipage}\hfill
  \begin{minipage}[c]{2.6cm}\raggedleft\includegraphics[width=2.6cm]{\fm@logo}\end{minipage}\par
  \vspace{2pt}}

% ---------- barres de secció del color del curs ----------
\renewcommand{\seccio}[2]{\par\needspace{\fmneed\baselineskip}\bigskip
  \begin{tcolorbox}[colback=curs,colframe=curs,coltext=curstext,arc=2mm,boxrule=0pt,left=3mm,right=3mm,top=1.2mm,bottom=1.2mm]
  {\large\bfseries #1\quad #2}
  \end{tcolorbox}\vspace{-2pt}}

% ---------- estrelles de dificultat del color del curs (a Batxillerat no se'n posen) ----------
\@ifundefined{estrelles}{}{%
  \renewcommand{\estrelles}[1]{{\color{curs}\ifcase#1\or\faStar\faStar[regular]\faStar[regular]\or\faStar\faStar\faStar[regular]\or\faStar\faStar\faStar\fi}}}

% ---------- capçalera i peu de les pàgines ----------
% primera pàgina: logos institucionals al peu
\fancypagestyle{primera}{\fancyhf{}\renewcommand{\headrulewidth}{0pt}
  \fancyfoot[L]{\raisebox{-0.32\height}{\includegraphics[height=0.78cm]{logo-gva.png}}\hspace{5mm}%
    \raisebox{-0.32\height}{\includegraphics[height=0.78cm]{logo-vermellar.png}}}
  \fancyfoot[R]{\scriptsize\color{muted}\href{https://creativecommons.org/licenses/by-nc-sa/4.0/deed.ca}{\ccbyncsa}\ \ CC BY-NC-SA 4.0 · fentmates.org\qquad\small\thepage}}
% la resta: el símbol a l'esquerra i el curs a la dreta, amb el títol al centre com abans
\fancyhead[L]{\raisebox{-0.3\height}{\includegraphics[height=0.62cm]{\fm@marca}}}
\fancyhead[R]{\small\bfseries\color{cursfosc}\fm@cursnom}
\renewcommand{\headrule}{\color{curs!45}\hrule height \headrulewidth}
\makeatother

%   \fmcurs{1eso}
% i, al començament del document:
%   \capcalera{Matemàtiques}{1r ESO}{Fitxa d'exercicis · Unitat 2}{Divisibilitat}
% Opcional abans de \input: \def\fmneed{14} (espai mínim abans de cada secció).
% =====================================================================
\makeatletter
\def\fm@logodir{../../logo/pdf}
\providecommand{\fmneed}{10}
\def\fm@cursnom{}

% clau del curs, color i color del text damunt del color
\newcommand{\fm@curs}[3]{\@namedef{fm@color@#1}{#2}\@namedef{fm@text@#1}{#3}}
\fm@curs{1eso}{E6B84A}{navy}
\fm@curs{2eso}{3B82F6}{white}
\fm@curs{3eso}{FF7A50}{white}
\fm@curs{4eso-a}{A78BFA}{white}
\fm@curs{4eso-b}{6D28D9}{white}
\fm@curs{1batx-mi}{9B1C1C}{white}
\fm@curs{1batx-ccss}{EF6B6B}{white}
\fm@curs{2batx-mii}{1E6B47}{white}
\fm@curs{2batx-ccss}{2BB3A3}{white}

\newcommand{\fmcurs}[1]{%
  \@ifundefined{fm@color@#1}{\PackageError{fentmates}{Curs desconegut: #1}{}}{}%
  \edef\fm@tmp{\noexpand\definecolor{curs}{HTML}{\@nameuse{fm@color@#1}}%
               \noexpand\colorlet{curstext}{\@nameuse{fm@text@#1}}}\fm@tmp
  \colorlet{cursfosc}{curs!72!black}%
  \def\fm@logo{\fm@logodir/vertical-#1.pdf}%
  \def\fm@marca{\fm@logodir/marca-#1.pdf}}

% ---------- capçalera de la primera pàgina ----------
% #1 matèria (petit), #2 curs, #3 tipus de document i unitat, #4 títol
\newcommand{\capcalera}[4]{%
  \gdef\fm@cursnom{#2}%
  \par\noindent
  \begin{tcolorbox}[colback=curs,colframe=curs,boxrule=0pt,arc=2.2mm,width=3.5cm,height=2.15cm,
      valign=center,left=3mm,right=2mm,top=1mm,bottom=1mm,nobeforeafter,box align=center]
    \color{curstext}{\footnotesize #1}\par\vspace{1pt}{\fontsize{19}{21}\selectfont\bfseries #2\par}%
  \end{tcolorbox}\hspace{5mm}%
  \begin{minipage}[c]{\dimexpr\linewidth-3.5cm-5mm-2.9cm\relax}
    {\color{cursfosc}\small\bfseries\MakeUppercase{#3}}\par\vspace{2pt}
    {\color{navy}\fontsize{20}{23}\selectfont\bfseries #4\par}
  \end{minipage}\hfill
  \begin{minipage}[c]{2.6cm}\raggedleft\includegraphics[width=2.6cm]{\fm@logo}\end{minipage}\par
  \vspace{2pt}}

% ---------- barres de secció del color del curs ----------
\renewcommand{\seccio}[2]{\par\needspace{\fmneed\baselineskip}\bigskip
  \begin{tcolorbox}[colback=curs,colframe=curs,coltext=curstext,arc=2mm,boxrule=0pt,left=3mm,right=3mm,top=1.2mm,bottom=1.2mm]
  {\large\bfseries #1\quad #2}
  \end{tcolorbox}\vspace{-2pt}}

% ---------- estrelles de dificultat del color del curs (a Batxillerat no se'n posen) ----------
\@ifundefined{estrelles}{}{%
  \renewcommand{\estrelles}[1]{{\color{curs}\ifcase#1\or\faStar\faStar[regular]\faStar[regular]\or\faStar\faStar\faStar[regular]\or\faStar\faStar\faStar\fi}}}

% ---------- capçalera i peu de les pàgines ----------
% primera pàgina: logos institucionals al peu
\fancypagestyle{primera}{\fancyhf{}\renewcommand{\headrulewidth}{0pt}
  \fancyfoot[L]{\raisebox{-0.32\height}{\includegraphics[height=0.78cm]{logo-gva.png}}\hspace{5mm}%
    \raisebox{-0.32\height}{\includegraphics[height=0.78cm]{logo-vermellar.png}}}
  \fancyfoot[R]{\scriptsize\color{muted}\href{https://creativecommons.org/licenses/by-nc-sa/4.0/deed.ca}{\ccbyncsa}\ \ CC BY-NC-SA 4.0 · fentmates.org\qquad\small\thepage}}
% la resta: el símbol a l'esquerra i el curs a la dreta, amb el títol al centre com abans
\fancyhead[L]{\raisebox{-0.3\height}{\includegraphics[height=0.62cm]{\fm@marca}}}
\fancyhead[R]{\small\bfseries\color{cursfosc}\fm@cursnom}
\renewcommand{\headrule}{\color{curs!45}\hrule height \headrulewidth}
\makeatother

%   \fmcurs{1eso}
% i, al començament del document:
%   \capcalera{Matemàtiques}{1r ESO}{Fitxa d'exercicis · Unitat 2}{Divisibilitat}
% Opcional abans de \input: \def\fmneed{14} (espai mínim abans de cada secció).
% =====================================================================
\makeatletter
\def\fm@logodir{../../logo/pdf}
\providecommand{\fmneed}{10}
\def\fm@cursnom{}

% clau del curs, color i color del text damunt del color
\newcommand{\fm@curs}[3]{\@namedef{fm@color@#1}{#2}\@namedef{fm@text@#1}{#3}}
\fm@curs{1eso}{E6B84A}{navy}
\fm@curs{2eso}{3B82F6}{white}
\fm@curs{3eso}{FF7A50}{white}
\fm@curs{4eso-a}{A78BFA}{white}
\fm@curs{4eso-b}{6D28D9}{white}
\fm@curs{1batx-mi}{9B1C1C}{white}
\fm@curs{1batx-ccss}{EF6B6B}{white}
\fm@curs{2batx-mii}{1E6B47}{white}
\fm@curs{2batx-ccss}{2BB3A3}{white}

\newcommand{\fmcurs}[1]{%
  \@ifundefined{fm@color@#1}{\PackageError{fentmates}{Curs desconegut: #1}{}}{}%
  \edef\fm@tmp{\noexpand\definecolor{curs}{HTML}{\@nameuse{fm@color@#1}}%
               \noexpand\colorlet{curstext}{\@nameuse{fm@text@#1}}}\fm@tmp
  \colorlet{cursfosc}{curs!72!black}%
  \def\fm@logo{\fm@logodir/vertical-#1.pdf}%
  \def\fm@marca{\fm@logodir/marca-#1.pdf}}

% ---------- capçalera de la primera pàgina ----------
% #1 matèria (petit), #2 curs, #3 tipus de document i unitat, #4 títol
\newcommand{\capcalera}[4]{%
  \gdef\fm@cursnom{#2}%
  \par\noindent
  \begin{tcolorbox}[colback=curs,colframe=curs,boxrule=0pt,arc=2.2mm,width=3.5cm,height=2.15cm,
      valign=center,left=3mm,right=2mm,top=1mm,bottom=1mm,nobeforeafter,box align=center]
    \color{curstext}{\footnotesize #1}\par\vspace{1pt}{\fontsize{19}{21}\selectfont\bfseries #2\par}%
  \end{tcolorbox}\hspace{5mm}%
  \begin{minipage}[c]{\dimexpr\linewidth-3.5cm-5mm-2.9cm\relax}
    {\color{cursfosc}\small\bfseries\MakeUppercase{#3}}\par\vspace{2pt}
    {\color{navy}\fontsize{20}{23}\selectfont\bfseries #4\par}
  \end{minipage}\hfill
  \begin{minipage}[c]{2.6cm}\raggedleft\includegraphics[width=2.6cm]{\fm@logo}\end{minipage}\par
  \vspace{2pt}}

% ---------- barres de secció del color del curs ----------
\renewcommand{\seccio}[2]{\par\needspace{\fmneed\baselineskip}\bigskip
  \begin{tcolorbox}[colback=curs,colframe=curs,coltext=curstext,arc=2mm,boxrule=0pt,left=3mm,right=3mm,top=1.2mm,bottom=1.2mm]
  {\large\bfseries #1\quad #2}
  \end{tcolorbox}\vspace{-2pt}}

% ---------- estrelles de dificultat del color del curs (a Batxillerat no se'n posen) ----------
\@ifundefined{estrelles}{}{%
  \renewcommand{\estrelles}[1]{{\color{curs}\ifcase#1\or\faStar\faStar[regular]\faStar[regular]\or\faStar\faStar\faStar[regular]\or\faStar\faStar\faStar\fi}}}

% ---------- capçalera i peu de les pàgines ----------
% primera pàgina: logos institucionals al peu
\fancypagestyle{primera}{\fancyhf{}\renewcommand{\headrulewidth}{0pt}
  \fancyfoot[L]{\raisebox{-0.32\height}{\includegraphics[height=0.78cm]{logo-gva.png}}\hspace{5mm}%
    \raisebox{-0.32\height}{\includegraphics[height=0.78cm]{logo-vermellar.png}}}
  \fancyfoot[R]{\scriptsize\color{muted}\href{https://creativecommons.org/licenses/by-nc-sa/4.0/deed.ca}{\ccbyncsa}\ \ CC BY-NC-SA 4.0 · fentmates.org\qquad\small\thepage}}
% la resta: el símbol a l'esquerra i el curs a la dreta, amb el títol al centre com abans
\fancyhead[L]{\raisebox{-0.3\height}{\includegraphics[height=0.62cm]{\fm@marca}}}
\fancyhead[R]{\small\bfseries\color{cursfosc}\fm@cursnom}
\renewcommand{\headrule}{\color{curs!45}\hrule height \headrulewidth}
\makeatother

\fmcurs{3eso}

\begin{document}
\thispagestyle{primera}

\capcalera{Matemáticas}{3\textsuperscript{o} ESO}{Ficha de ejercicios · Unidad 9}{Perímetros y áreas. Teorema de Pitágoras}
\vspace{10pt}
Nombre y apellidos: \rule{6.6cm}{0.4pt}\hfill Grupo: \rule{1.6cm}{0.4pt}\hfill Fecha: \rule{2.2cm}{0.4pt}
\vspace{8pt}

\begin{tcolorbox}[colback=soft,colframe=soft,arc=2mm,left=3mm,right=3mm,top=1.5mm,bottom=1.5mm]
\small Ejercicios graduados por niveles de dificultad:\quad \estrelles{1}\ básico \qquad \estrelles{2}\ medio \qquad \estrelles{3}\ avanzado.\\[2pt] Resuélvelos en el cuaderno mostrando todos los pasos y haz un dibujo de cada situación. Si un resultado no es exacto, redondéalo a las centésimas; para $\pi$, usa la tecla de la calculadora.
\end{tcolorbox}

% =====================================================
\seccio{1}{El teorema de Pitágoras}

\ex{1} Calcula la hipotenusa de un triángulo rectángulo cuyos catetos miden (si no es exacta, redondea a las centésimas):\begin{tasks}(3)
  \task 6 cm y 8 cm
  \task 5 cm y 12 cm
  \task 9 m y 40 m
  \task 3 cm y 7 cm
  \task 1,5 m y 2 m
  \task 10 cm y 10 cm
\end{tasks}\recorda{En un triángulo rectángulo, $c^2=a^2+b^2$, donde $c$ es la \textbf{hipotenusa} (el lado opuesto al ángulo recto, el más largo) y $a$, $b$ son los \textbf{catetos}. Para la hipotenusa sumamos: $c=\sqrt{a^2+b^2}$; para un cateto restamos: $a=\sqrt{c^2-b^2}$.}

\ex{1} Calcula el cateto que falta en cada triángulo rectángulo ($c$ es la hipotenusa):\begin{tasks}(3)
  \task $c=25$, $b=7$
  \task $c=17$, $a=8$
  \task $c=10$, $b=8$
  \task $c=12$, $b=5$
  \task $c=6{,}5$, $a=2{,}5$
  \task $c=9$, $a=4$
\end{tasks}

\ex{1} Calcula la longitud $x$ en cada triángulo rectángulo (medidas en cm):\par\smallskip\begin{center}\begin{tikzpicture}[x=0.15cm,y=0.15cm,baseline=(current bounding box.center),line join=round]
\draw[fill=accent!22,draw=accent,thick] (0.000,0.000) -- (12.000,0.000) -- (0.000,16.000) -- cycle;
\draw[thin] (1.467,0.000) -- (1.467,1.467) -- (0.000,1.467);
\node[font=\small,below] at (6.000,0.000) {12};
\node[font=\small,left] at (0.000,8.000) {16};
\node[font=\small,above right] at (6.000,8.000) {$x$};
\end{tikzpicture}\qquad \begin{tikzpicture}[x=0.15cm,y=0.15cm,baseline=(current bounding box.center),line join=round]
\draw[fill=accent!22,draw=accent,thick] (0.000,0.000) -- (12.000,0.000) -- (0.000,5.000) -- cycle;
\draw[thin] (1.467,0.000) -- (1.467,1.467) -- (0.000,1.467);
\node[font=\small,below] at (6.000,0.000) {$x$};
\node[font=\small,left] at (0.000,2.500) {5};
\node[font=\small,above right] at (6.000,2.500) {13};
\end{tikzpicture}\qquad \begin{tikzpicture}[x=0.15cm,y=0.15cm,baseline=(current bounding box.center),line join=round]
\draw[fill=accent!22,draw=accent,thick] (0.000,0.000) -- (15.000,0.000) -- (0.000,8.000) -- cycle;
\draw[thin] (1.467,0.000) -- (1.467,1.467) -- (0.000,1.467);
\node[font=\small,below] at (7.500,0.000) {15};
\node[font=\small,left] at (0.000,4.000) {$x$};
\node[font=\small,above right] at (7.500,4.000) {17};
\end{tikzpicture}\end{center}

\ex{2} Una escalera de 6,5 m está apoyada en una pared. El pie de la escalera está a 2,5 m de la pared. ¿A qué altura de la pared llega?

\ex{2} Volamos una cometa con un hilo de 50 m tenso. Si estamos a 30 m (en horizontal) del punto del suelo que hay justo debajo de la cometa, ¿a qué altura vuela? (No tengas en cuenta la altura de la mano.)

\ex{2} Un barco sale del puerto y navega 9 km hacia el norte y, después, 12 km hacia el este. ¿A qué distancia en línea recta se encuentra del puerto?

\ex{2} Una rampa salva un desnivel de 0,8 m y ocupa 6 m en horizontal. ¿Cuánto mide la rampa?

\ex{3} En un triángulo rectángulo, un cateto mide 2 cm más que el otro y la hipotenusa mide 10 cm. ¿Cuánto miden los catetos?

\ex{3} Un poste vertical de 9 m se rompe por el viento y la punta toca el suelo a 3 m de la base (la parte de abajo sigue recta). ¿A qué altura se ha roto?


% =====================================================
\seccio{2}{Clasificación de triángulos por los lados}

\ex{1} Clasifica según los ángulos (acutángulo, rectángulo u obtusángulo) los triángulos de lados:\begin{tasks}(3)
  \task 5, 12 y 13
  \task 6, 7 y 8
  \task 4, 6 y 9
  \task 10, 24 y 26
  \task 7, 8 y 12
  \task 9, 10 y 11
\end{tasks}\recorda{Si $c$ es el lado más largo: $c^2<a^2+b^2$ \ $\Rightarrow$ acutángulo; \ $c^2=a^2+b^2$ \ $\Rightarrow$ rectángulo; \ $c^2>a^2+b^2$ \ $\Rightarrow$ obtusángulo.}

\ex{1} ¿Cuáles de estas ternas son pitagóricas?\begin{tasks}(5)
  \task 12, 16, 20
  \task 9, 12, 16
  \task 15, 36, 39
  \task 11, 60, 61
  \task 8, 9, 12
\end{tasks}

\ex{2} ¿Pueden ser los lados de un triángulo? Si es así, clasifícalo.\begin{tasks}(3)
  \task 3, 5 y 9
  \task 4, 4 y 7
  \task 5, 7 y 12
\end{tasks}

\ex{2} Un triángulo tiene los lados de 1,5 m, 2 m y 2,5 m. ¿Es rectángulo? Si lo es, calcula su área tomando los catetos como base y altura.

\ex{2} Un carpintero quiere comprobar si una estantería de 90 cm × 120 cm tiene los ángulos rectos. ¿Cuánto tiene que medir la diagonal? Si la mide y le da 152 cm, ¿cómo es el ángulo?

\ex{3} ¿Cuánto tiene que medir $x$ para que el triángulo de lados 9, 12 y $x$ sea rectángulo? Hay dos posibilidades.

\ex{3} El triángulo tiene lados 6, 8 y $x$, con $x$ un número entero mayor que 8. ¿Para qué valores de $x$ es acutángulo? ¿Y rectángulo? ¿Y obtusángulo?


% =====================================================
\seccio{3}{Aplicaciones del teorema de Pitágoras}

\ex{1} Calcula la diagonal de:\begin{tasks}(3)
  \task un rectángulo de 20 × 15 cm
  \task un cuadrado de 7 cm de lado
  \task una puerta de 2,4 × 1 m
\end{tasks}\recorda{Busca el triángulo rectángulo escondido: la diagonal de un rectángulo, la altura de un triángulo isósceles o equilátero, la mitad de las diagonales de un rombo o la altura de un trapecio forman uno. Decide qué lado es la hipotenusa y aplica Pitágoras.}

\ex{1} Calcula la altura de un triángulo isósceles de lados iguales de 17 cm y base de 16 cm, y la de un triángulo equilátero de 12 cm de lado.

\ex{1} Las diagonales de un rombo miden 18 cm y 24 cm. Calcula su lado y su perímetro.

\ex{2} La diagonal de un cuadrado mide 12 cm. Calcula su lado y su perímetro.

\ex{2} Un trapecio rectángulo tiene las bases de 15 cm y 9 cm y la altura de 8 cm. Calcula el lado oblicuo y el perímetro.

\ex{2} Un trapecio isósceles tiene las bases de 22 cm y 12 cm y los lados no paralelos de 13 cm. Calcula su altura.

\ex{2} Calcula la distancia entre los puntos:\begin{tasks}(3)
  \task $A(2,\,1)$ y $B(5,\,5)$
  \task $C(-3,\,4)$ y $D(3,\,-4)$
  \task $E(0,\,-2)$ y $F(4,\,3)$
\end{tasks}\recorda{$d(A,\,B)=\sqrt{(x_2-x_1)^2+(y_2-y_1)^2}$: es la hipotenusa de un triángulo rectángulo cuyos catetos son la diferencia de las $x$ y la de las $y$.}

\ex{3} Sea el triángulo de vértices $A(1,\,1)$, $B(7,\,1)$ y $C(4,\,5)$. Calcula sus lados y su perímetro, y di cómo es según los lados y según los ángulos.

\ex{3} Un rombo tiene 10 cm de lado y una diagonal de 12 cm. ¿Cuánto mide la otra diagonal?

\ex{3} Un televisor de 32 pulgadas tiene la pantalla de proporción 16:9 (la anchura y la altura son $16k$ y $9k$). Las pulgadas miden la diagonal y 1 pulgada son 2,54 cm. Calcula la anchura y la altura de la pantalla.


% =====================================================
\seccio{4}{Perímetros y áreas de triángulos y cuadriláteros}

\ex{1} Calcula el área y el perímetro de:\begin{tasks}(2)
  \task un cuadrado de 9 cm de lado
  \task un rectángulo de 12 cm × 5 cm
  \task un romboide de base 10 cm, altura 6 cm y lado oblicuo 7 cm
  \task un rombo de diagonales 16 cm y 12 cm
\end{tasks}\recorda{Cuadrado $A=l^2$; \ rectángulo y romboide $A=b\cdot h$; \ rombo $A=\frac{D\cdot d}{2}$; \ trapecio $A=\frac{(B+b)\cdot h}{2}$; \ triángulo $A=\frac{b\cdot h}{2}$. El perímetro es la suma de los lados.}

\ex{1} Calcula el área de:\begin{tasks}(2)
  \task un trapecio de bases 11 m y 7 m y altura 5 m
  \task un triángulo de 14 cm de base y 9 cm de altura
\end{tasks}

\ex{1} Expresa en la unidad que se indica:\begin{tasks}(4)
  \task $3{,}5\text{ m}^2$ en cm²
  \task $45\,000\text{ cm}^2$ en m²
  \task $2{,}3\text{ ha}$ en m²
  \task $750\text{ dm}^2$ en m²
\end{tasks}

\ex{1} Un rectángulo tiene 72 cm² de área y 9 cm de base. Calcula su altura y su perímetro.

\ex{2} Calcula el área de un triángulo equilátero de 8 cm de lado.

\ex{2} Calcula el área y el perímetro de un triángulo isósceles de lados iguales de 10 cm y base de 12 cm.

\ex{2} Calcula el área y el perímetro de un trapecio isósceles de bases 20 cm y 10 cm y lados no paralelos de 13 cm.

\ex{2} Un rombo tiene 40 cm de perímetro y una diagonal de 16 cm. Calcula su área.

\ex{2} Un campo rectangular mide 120 m × 85 m. Expresa su superficie en m² y en hectáreas. ¿Cuánto costará sembrar césped a 3 € el m²?

\ex{3} Queremos pintar una pared de 4,5 m × 2,6 m que tiene una ventana de 1,2 m × 1 m. Si un bote de pintura cubre 4 m², ¿cuántos botes tenemos que comprar?

\ex{3} Los catetos de un triángulo rectángulo miden 20 cm y 15 cm. Calcula su área, la hipotenusa y la altura sobre la hipotenusa.

\ex{3} Un rectángulo tiene 34 cm de perímetro y 13 cm de diagonal. ¿Cuánto miden los lados? ¿Cuál es el área?


% =====================================================
\seccio{5}{Polígonos regulares}

\ex{1} Calcula el perímetro y el área de los polígonos regulares:\begin{tasks}(1)
  \task pentágono de 6 cm de lado y 4,13 cm de apotema
  \task hexágono de 5 cm de lado y 4,33 cm de apotema
  \task octógono de 4 cm de lado y 4,83 cm de apotema
\end{tasks}\recorda{En un polígono regular, $A=\frac{P\cdot a}{2}$, donde $P$ es el perímetro y $a$ la apotema. El radio, la apotema y medio lado forman un triángulo rectángulo: $r^2=a^2+\left(\frac l2\right)^2$. En el hexágono regular, $r=l$.}

\ex{1} Un hexágono regular tiene 8 cm de lado. ¿Cuánto mide el radio? Calcula su apotema y su área.

\ex{2} Un pentágono regular tiene 12 cm de lado y 10,2 cm de radio. Calcula su apotema y su área.

\ex{2} Un cuadrado está inscrito en una circunferencia de 6 cm de radio. Calcula su lado, su apotema y su área.

\ex{2} Queremos cubrir 3 m² de pared con baldosas en forma de hexágono regular de 10 cm de lado. ¿Cuántas baldosas hacen falta, como mínimo?

\ex{3} Un triángulo equilátero está inscrito en una circunferencia de 4 cm de radio. Sabiendo que su apotema es la mitad del radio, calcula su lado y su área.


% =====================================================
\seccio{6}{La circunferencia: longitud y arcos}

\ex{1} Calcula la longitud de la circunferencia de:\begin{tasks}(3)
  \task 7 cm de radio
  \task 15 cm de diámetro
  \task 2,5 m de radio
\end{tasks}\recorda{$L=2\pi r=\pi d$. \ Un arco de $n^\circ$ mide $L_{\text{arco}}=\frac{2\pi r\cdot n}{360}$. Usa la tecla $\pi$ de la calculadora y redondea solo al final.}

\ex{1} Una circunferencia mide 62,83 cm de longitud. ¿Cuánto miden el radio y el diámetro?

\ex{1} Calcula la longitud de los arcos:\begin{tasks}(3)
  \task de 90° y 8 cm de radio
  \task de 30° y 12 cm de radio
  \task de 225° y 4 cm de radio
\end{tasks}

\ex{2} Las ruedas de una bicicleta miden 60 cm de diámetro. ¿Cuántas vueltas da cada rueda en un recorrido de 1,5 km?

\ex{2} La cuerda interior de una pista de atletismo está formada por dos rectas de 100 m y dos semicircunferencias de 36,5 m de radio. ¿Cuánto mide una vuelta?

\ex{2} La aguja minutera de un reloj mide 15 cm. ¿Qué distancia recorre su punta en 20 minutos?

\ex{3} En una circunferencia de 12 cm de radio, un arco mide 18,85 cm. ¿Cuántos grados tiene?

\ex{3} La Tierra da una vuelta al Sol cada año siguiendo una órbita casi circular de 150 millones de km de radio. ¿A qué velocidad media (en km/h) viaja la Tierra?


% =====================================================
\seccio{7}{Área del círculo y figuras circulares}

\ex{1} Calcula el área del círculo de:\begin{tasks}(3)
  \task 5 cm de radio
  \task 12 cm de diámetro
  \task 0,8 m de radio
\end{tasks}\recorda{Círculo $A=\pi r^2$; \ sector de $n^\circ$: $A=\frac{\pi r^2\cdot n}{360}$; \ corona: $A=\pi\left(R^2-r^2\right)$; \ trapecio circular: $A=\frac{\pi\left(R^2-r^2\right)\cdot n}{360}$.}

\ex{1} Calcula el área de los sectores circulares:\begin{tasks}(3)
  \task de 60° y 6 cm de radio
  \task de 150° y 10 cm de radio
  \task de 270° y 2 cm de radio
\end{tasks}

\ex{1} Calcula el área de una corona circular de radios 9 cm y 4 cm.

\ex{2} Un círculo tiene 200 cm² de área. Calcula su radio y la longitud de la circunferencia.

\ex{2} Calcula el área de un trapecio circular de radios 12 cm y 7 cm y amplitud 80°.

\ex{2} Una pizza de 30 cm de diámetro se corta en 8 porciones iguales. Calcula la superficie de cada porción y la longitud de su borde (el arco).

\ex{2} Un hexágono regular tiene 30 cm de perímetro. Calcula el área del círculo circunscrito (el que pasa por los vértices) y la del círculo inscrito (el radio es la apotema).

\ex{3} Calcula el área del segmento circular sombreado, determinado por un ángulo de 60° en un círculo de 6 cm de radio. (El triángulo es equilátero.)\par\smallskip\begin{center}\begin{tikzpicture}[x=0.55cm,y=0.55cm,baseline=(current bounding box.center),line join=round]
\draw[fill=navy!18,draw=navy,thick] (0.000,0.000) -- (3.000,0.000) arc[start angle=0,end angle=60,radius=3] -- cycle;
\draw[fill=white,draw=white] (0.000,0.000) -- (3.000,0.000) -- (1.500,2.598) -- cycle;
\draw[thick] (0.000,0.000) -- (3.000,0.000);
\draw[thick] (0.000,0.000) -- (1.500,2.598);
\draw[thick] (3.000,0.000) -- (1.500,2.598);
\draw[thick] (3.000,0.000) arc[start angle=0,end angle=60,radius=3];
\fill (0.000,0.000) circle (1.4pt);
\node[font=\small,below] at (1.500,0.000) {6 cm};
\draw[thick] (0.600,0.000) arc[start angle=0,end angle=60,radius=0.6];
\node[font=\small,right] at (0.750,0.350) {\scriptsize 60°};
\end{tikzpicture}\end{center}

\ex{3} Una corona circular tiene un área de $50\pi$ cm² y el radio exterior mide 10 cm. ¿Cuánto mide el radio interior?


% =====================================================
\seccio{8}{Figuras compuestas y problemas}

\ex{1} Calcula el área y el perímetro de la fachada de la caseta.\par\smallskip\begin{center}\begin{tikzpicture}[x=0.17cm,y=0.17cm,baseline=(current bounding box.center),line join=round]
\draw[fill=teal!12,draw=teal,thick] (0.000,0.000) -- (6.000,0.000) -- (6.000,6.000) -- (3.000,10.000) -- (0.000,6.000) -- cycle;
\draw[dashed,muted] (0.000,6.000) -- (6.000,6.000);
\draw[dashed,muted] (3.000,6.000) -- (3.000,10.000);
\draw[thin] (4.294,6.000) -- (4.294,7.294) -- (3.000,7.294);
\node[font=\small,below] at (3.000,0.000) {6 m};
\node[font=\small,right] at (6.000,3.000) {6 m};
\node[font=\small,right] at (3.000,8.000) {4 m};
\end{tikzpicture}\end{center}\recorda{Para una figura compuesta, descomponla en figuras conocidas y suma las áreas, o mírala como una figura grande a la que le quitas trozos y resta. El perímetro es solo el contorno exterior.}

\ex{1} A un cuadrado de 8 cm de lado le añadimos un semicírculo sobre uno de los lados. Calcula el área y el perímetro de la figura.\par\smallskip\begin{center}\begin{tikzpicture}[x=0.15cm,y=0.15cm,baseline=(current bounding box.center),line join=round]
\draw[fill=navy!18,draw=navy,thick] (0.000,0.000) -- (8.000,0.000) -- (8.000,8.000) -- (0.000,8.000) -- cycle;
\draw[fill=navy!18,draw=navy,thick] (4.000,8.000) -- (8.000,8.000) arc[start angle=0,end angle=180,radius=4] -- cycle;
\draw[dashed,muted] (0.000,8.000) -- (8.000,8.000);
\node[font=\small,below] at (4.000,0.000) {8 cm};
\end{tikzpicture}\end{center}

\ex{2} De un rectángulo de 10 cm × 6 cm recortamos dos semicírculos de 6 cm de diámetro en los lados cortos. Calcula el área y el perímetro de la figura que queda.\par\smallskip\begin{center}\begin{tikzpicture}[x=0.2cm,y=0.2cm,baseline=(current bounding box.center),line join=round]
\draw[fill=navy!18,draw=navy,thick] (0.000,0.000) -- (10.000,0.000) -- (10.000,6.000) -- (0.000,6.000) -- cycle;
\draw[fill=white,draw=white] (0.000,3.000) -- (0.000,0.000) arc[start angle=-90,end angle=90,radius=3] -- cycle;
\draw[fill=white,draw=white] (10.000,3.000) -- (10.000,6.000) arc[start angle=90,end angle=270,radius=3] -- cycle;
\draw[thick] (0.000,0.000) arc[start angle=-90,end angle=90,radius=3];
\draw[thick] (10.000,6.000) arc[start angle=90,end angle=270,radius=3];
\node[font=\small,below] at (5.000,0.000) {10 cm};
\node[font=\small,right] at (10.000,3.000) {6 cm};
\end{tikzpicture}\end{center}

\ex{2} Calcula el área de la zona sombreada: un cuadrado de 12 cm de lado menos un cuarto de círculo de radio 12 cm.\par\smallskip\begin{center}\begin{tikzpicture}[x=0.13cm,y=0.13cm,baseline=(current bounding box.center),line join=round]
\draw[fill=navy!18,draw=navy,thick] (0.000,0.000) -- (12.000,0.000) -- (12.000,12.000) -- (0.000,12.000) -- cycle;
\draw[fill=white,draw=white] (0.000,0.000) -- (12.000,0.000) arc[start angle=0,end angle=90,radius=12] -- cycle;
\draw[thick] (12.000,0.000) arc[start angle=0,end angle=90,radius=12];
\node[font=\small,below] at (6.000,0.000) {12 cm};
\draw[dashed,muted] (0.000,0.000) -- (8.485,8.485);
\node[font=\small,above left] at (4.200,4.200) {12};
\end{tikzpicture}\end{center}

\ex{2} Una ventana tiene forma de rectángulo de 1,4 m de anchura y 1,8 m de altura terminado en un semicírculo. Calcula la superficie de cristal y la longitud del marco exterior.\par\smallskip\begin{center}\begin{tikzpicture}[x=0.75cm,y=0.75cm,baseline=(current bounding box.center),line join=round]
\draw[fill=teal!12,draw=teal,thick] (0.000,0.000) -- (1.400,0.000) -- (1.400,1.800) -- (0.000,1.800) -- cycle;
\draw[fill=teal!12,draw=teal,thick] (0.700,1.800) -- (1.400,1.800) arc[start angle=0,end angle=180,radius=0.7] -- cycle;
\draw[dashed,muted] (0.000,1.800) -- (1.400,1.800);
\node[font=\small,below] at (0.700,0.000) {1,4 m};
\node[font=\small,right] at (1.400,0.900) {1,8 m};
\end{tikzpicture}\end{center}

\ex{2} Alrededor de una piscina de 25 m × 12 m se hace una acera de 2 m de anchura. Calcula la superficie de la acera.

\ex{3} Calcula el área de la «hoja» sombreada. Los arcos son cuartos de circunferencia con centro en dos vértices opuestos del cuadrado de 10 cm de lado.\par\smallskip\begin{center}\begin{tikzpicture}[x=0.15cm,y=0.15cm,baseline=(current bounding box.center),line join=round]
\draw[fill=teal!12,draw=teal,thick] (0.000,0.000) -- (10.000,0.000) -- (10.000,10.000) -- (0.000,10.000) -- cycle;
\draw[fill=navy!18,draw=navy,thick] (10,0) arc[start angle=0,end angle=90,radius=10] arc[start angle=180,end angle=270,radius=10] -- cycle;
\node[font=\small,below] at (5.000,0.000) {10 cm};
\end{tikzpicture}\end{center}

\ex{3} Tres círculos iguales caben justo dentro de un rectángulo de 30 cm × 10 cm. Calcula el área de la zona sombreada.\par\smallskip\begin{center}\begin{tikzpicture}[x=0.14cm,y=0.14cm,baseline=(current bounding box.center),line join=round]
\draw[fill=navy!18,draw=navy,thick] (0.000,0.000) -- (30.000,0.000) -- (30.000,10.000) -- (0.000,10.000) -- cycle;
\draw[fill=white,draw=white] (5.000,5.000) circle[radius=5];
\draw[fill=white,draw=white] (15.000,5.000) circle[radius=5];
\draw[fill=white,draw=white] (25.000,5.000) circle[radius=5];
\draw[draw=teal,thick] (5.000,5.000) circle[radius=5];
\draw[draw=teal,thick] (15.000,5.000) circle[radius=5];
\draw[draw=teal,thick] (25.000,5.000) circle[radius=5];
\draw[dashed,muted] (0.000,0.000) -- (30.000,0.000) -- (30.000,10.000) -- (0.000,10.000) -- cycle;
\node[font=\small,below] at (15.000,0.000) {30 cm};
\node[font=\small,right] at (30.000,5.000) {10 cm};
\end{tikzpicture}\end{center}

\ex{3} Calcula el área comprendida entre un hexágono regular de 4 cm de lado y la circunferencia que pasa por sus vértices.

\ex{3} Una cabra está atada con una cuerda de 5 m a una esquina exterior de un corral cuadrado de 4 m de lado, y no puede entrar en el corral. ¿Qué superficie puede pastar?\par\smallskip\begin{center}\begin{tikzpicture}[x=0.24cm,y=0.24cm,baseline=(current bounding box.center),line join=round]
\draw[fill=navy!18,draw=navy,thick] (0.000,0.000) -- (0.000,5.000) arc[start angle=90,end angle=360,radius=5] -- cycle;
\draw[fill=navy!18,draw=navy,thick] (4.000,0.000) -- (5.000,0.000) arc[start angle=0,end angle=90,radius=1] -- cycle;
\draw[fill=navy!18,draw=navy,thick] (0.000,4.000) -- (1.000,4.000) arc[start angle=0,end angle=90,radius=1] -- cycle;
\draw[fill=teal!12,draw=teal,thick] (0.000,0.000) -- (4.000,0.000) -- (4.000,4.000) -- (0.000,4.000) -- cycle;
\fill (0.000,0.000) circle (1.4pt);
\draw[dashed,muted] (0.000,0.000) -- (-3.540,-3.540);
\node[font=\small,above left] at (-1.800,-1.800) {5 m};
\node[font=\small,inner sep=1pt] at (2.000,2.000) {corral};
\node[font=\small,below] at (2.000,4.000) {4 m};
\end{tikzpicture}\end{center}

\vfill
\begin{center}\color{muted}\small\faCheckCircle[regular]\ \ ¡Ánimo! Si te cuesta algún ejercicio, repasa los apuntes interactivos de la unidad.\end{center}

\end{document}
